\documentclass{article}

% Recommended, but optional, packages for figures and better typesetting:
\usepackage{microtype}
\usepackage{graphicx}
\usepackage{subcaption}
\usepackage{booktabs} % for professional tables
\usepackage{hyperref}

% Attempt to make hyperref and algorithmic work together better:
\newcommand{\theHalgorithm}{\arabic{algorithm}}

\usepackage{icml2026}

\renewcommand{\topfraction}{0.95}
\renewcommand{\bottomfraction}{0.95}
\renewcommand{\textfraction}{0.05}
\renewcommand{\floatpagefraction}{0.9}
\renewcommand{\dbltopfraction}{0.95}
\renewcommand{\dblfloatpagefraction}{0.9}

\setlength{\textfloatsep}{6pt plus 2pt minus 4pt}
\setlength{\floatsep}{6pt plus 2pt minus 2pt}
\setlength{\intextsep}{6pt plus 2pt minus 2pt}
\setlength{\dbltextfloatsep}{6pt plus 2pt minus 4pt}
\setlength{\dblfloatsep}{6pt plus 2pt minus 2pt}

\usepackage{amsmath}
\usepackage{amssymb}
\usepackage{mathtools}
\usepackage{amsthm}
\usepackage{algorithm}
\usepackage{algorithmic}
\usepackage{float}

% if you use cleveref..
\usepackage[capitalize,noabbrev]{cleveref}

%%%%%%%%%%%%%%%%%%%%%%%%%%%%%%%%
% THEOREMS
%%%%%%%%%%%%%%%%%%%%%%%%%%%%%%%%
\theoremstyle{plain}
\newtheorem{theorem}{Theorem}[section]
\newtheorem{proposition}[theorem]{Proposition}
\newtheorem{lemma}[theorem]{Lemma}
\newtheorem{corollary}[theorem]{Corollary}
\theoremstyle{definition}
\newtheorem{definition}[theorem]{Definition}
\newtheorem{assumption}[theorem]{Assumption}
\theoremstyle{remark}
\newtheorem{remark}[theorem]{Remark}

\icmltitlerunning{Data-Free Bayesian Test-Time Model Merging}

\begin{document}

\twocolumn[
  \icmltitle{Data-Free Bayesian Test-Time Model Merging with \\
    Soft Batch Normalization Buffer Fusion}

  \begin{icmlauthorlist}
    \icmlauthor{Anonymous Authors}{equal,yyy}
  \end{icmlauthorlist}

  \icmlaffiliation{yyy}{Department of Computer Science, Anonymous University, Location, Country}
  \icmlcorrespondingauthor{Anonymous Authors}{anon.email@domain.com}

  \icmlkeywords{Machine Learning, Model Merging, Test-Time Adaptation, Bayesian Inference}

  \vskip 0.3in
]

\printAffiliationsAndNotice{}

\begin{abstract}
Test-Time Model Merging (TTMM) is an emerging paradigm for dynamically fusing multiple
task-specific expert neural networks into a single multi-task model at deployment time. It adapts to
a non-stationary stream of test data without parameter backpropagation. However, existing TTMM
methods face three fundamental challenges: (1) a critical reliance on private source calibration
data for parameter-sensitivity (Fisher Information) estimation, (2) severe activation mismatches due
to the omission or hard-switching of Batch Normalization (BN) running statistics, and (3) a
susceptibility to the ``feedback trap'' (a failure mode of unconstrained entropy minimization
where parameters irreversibly collapse onto a single confidently wrong expert, which is particularly
severe when experts share a classification head). To resolve
these limitations, we propose \textbf{DF-Bayes-TTMM} (Data-Free Bayesian Test-Time Model Merging
with Soft BN Buffer Fusion). First, we formulate TTMM as a dynamic Bayesian mixture-of-experts,
estimating soft, continuous posterior probabilities for each active expert on-the-fly based on
predictive confidence. Second, we introduce Soft BN Buffer Fusion, which continuously blends expert
running statistics using the posterior weights to preserve activation integrity without
hard-switching discontinuities. We mathematically prove that our fusion equations perfectly
reconstruct the moments of the underlying mixture-of-Gaussians activation distribution. Third, we
establish a fully data-free, uncertainty-guided novelty detection mechanism based on the average
expert entropy. Our extensive empirical evaluations on non-stationary vision streams show that
DF-Bayes-TTMM completely eliminates private source data dependency, avoids the feedback trap, and
achieves state-of-the-art performance on open-world benchmarks while demonstrating competitive performance
and unique robustness trade-offs on deeper, more complex architectures. In
doing so, it overcomes the catastrophic collapse of unconstrained entropy minimization (improving
accuracy by up to 28\%) and outperforms competitive hard-routing baselines in open-world settings
while maintaining a continuous and principled design.
\end{abstract}

\section{Introduction}\vspace{-0.1in}
\label{sec:introduction}

The paradigm of pre-training large foundation models and fine-tuning them on specialized downstream
domains has led to a proliferation of specialized ``expert'' models \cite{Wortsman2022,
Radford2021}. While each expert excels on its specific task, deploying multiple independent experts
for multi-task applications is computationally prohibitive. To integrate disparate expert
capabilities at inference time, model merging techniques such as task arithmetic \cite{Ilharco2023},
Ties-Merging \cite{Yadav2023}, and Fisher weighted averaging \cite{Matena2022} have been developed.
These methods interpolate the weights of specialized models in parameter space to yield a single
unified model without joint training.

To handle dynamic shifts, Test-Time Model Merging (TTMM) adapts merging coefficients on-the-fly via
unsupervised objectives \cite{Yang2023, Yang2024}. By tuning low-dimensional coefficients instead of
full parameters, TTMM offers a highly efficient alternative to standard Test-Time Adaptation (TTA)
\cite{Wang2021}, preventing representational collapse and forgetting with minimal compute.

Despite their potential, existing TTMM frameworks are restricted by three fundamental bottlenecks:
\begin{enumerate}
\item \textbf{Source Data Inaccessibility:} Existing methods pre-compute Fisher Information on
source datasets to scale learning rates. However, training data is often private or inaccessible.
Fully data-free preconditioning is essential.
\item \textbf{Activation Mismatch:} Existing TTMM focus solely on weight interpolation, omitting
Batch Normalization (BN) statistics. This causes severe activation mismatches. DR-Fisher's hard
switching introduces discontinuities.
\item \textbf{The Feedback Trap:} Unconstrained entropy minimization can cause coefficients to
collapse toward a single out-of-distribution (OOD) expert that is spuriously confident, a failure
mode that is particularly severe when experts share a classification head.
\end{enumerate}

To overcome these boundaries, we introduce \textbf{DF-Bayes-TTMM} (Data-Free Bayesian Test-Time
Model Merging with Soft BN Buffer Fusion), a mathematically rigorous, fully unsupervised, and
completely data-free TTMM framework.

We formulate TTMM as a dynamic Bayesian mixture-of-experts. Rather than performing hard routing, we
estimate the continuous posterior probabilities of the experts on-the-fly using their predictive
confidence on the incoming test batch. This posterior distribution provides a principled,
soft-routing weight vector that guides both parameter and buffer fusion.

To resolve activation mismatches, we propose \textbf{Soft BN Buffer Fusion}, dynamically blending
expert running statistics using posterior weights. This blends moments of the underlying
mixture-of-Gaussians distribution, preserving activation integrity without discontinuities.

Furthermore, we design an uncertainty-guided novelty detection mechanism based on the average
expert entropy (which serves as a robust indicator of epistemic uncertainty). Our model robustly
identifies novel domains in the open-world stream without relying on private source prototypes or
distance thresholds.

Finally, we introduce a Maximum A Posteriori (MAP) adaptation step preconditioned by
\textbf{Test-Time Fisher (TT-Fisher)} sensitivities estimated on test samples. A temporal Gaussian
prior regularizes coefficient updates, ensuring stable, data-free adaptation.

We evaluate DF-Bayes-TTMM on non-stationary streams comprising MNIST, KMNIST, and FashionMNIST. Our
results show that DF-Bayes-TTMM completely avoids the feedback trap, achieves perfect novelty
routing, and delivers state-of-the-art performance, enhancing robustness while ensuring continuous,
data-free adaptation.

\section{Related Work}\vspace{-0.1in}
\label{sec:related_work}

\textbf{Model Merging:} Model merging aims to combine multiple neural networks trained on different
tasks or domains into a single unified network without the cost of retraining \cite{Wortsman2022,
Ainsworth2023, Jordan2022}. Standard weight averaging, however, suffers from parameter interference.
To mitigate this, Task Arithmetic \cite{Ilharco2023}, Ties-Merging \cite{Yadav2023}, and DARE
\cite{Yu2023} sparsify and scale weight deltas. Fisher Merging \cite{Matena2022} uses diagonal
Fisher Information matrices as parameter-wise weights. RegMean \cite{Jin2023} solves a quadratic
optimization problem to minimize activation distance. BECAME \cite{Li2025} proposes Bayesian
continual learning to compute closed-form merging coefficients. However, these methods are static
and require offline calibration data.

\textbf{Test-Time Adaptation (TTA):} TTA adapts a pre-trained model to target domain shifts
on-the-fly using unlabeled test streams \cite{Liang2023, Sun2020}. The foundational baseline Tent
\cite{Wang2021} optimizes scale and shift parameters of BN layers by minimizing prediction entropy.
Under continuous, non-i.i.d. shifts, TENT suffers from catastrophic parameter drift and
representational collapse. Advanced TTA methods use weight regularization \cite{Niu2022}, mean
teacher frameworks \cite{Wang2022}, and robust statistics estimation \cite{Yuan2023}. While
effective, full TTA is computationally expensive and prone to collapse under severe shifts.

\textbf{Test-Time Model Merging (TTMM):} TTMM combines the parameter efficiency of model merging
with the dynamic flexibility of TTA by adapting merging coefficients online \cite{Yang2024,
Zhao2024}. AdaMerging \cite{Yang2023, Yang2024} adaptively learns task-wise or layer-wise merging
coefficients by minimizing prediction entropy with a uniform learning rate, which often leads to
representational collapse in sensitive layers. FP-CA \cite{Luan2026} and IGGS-Merge
\cite{IGGSMerge2026} use pre-computed Fisher Information to precondition and damp updates in
sensitive layers. Recently, PROTO-TTMM \cite{ProtoTTMM2026} addresses open-world settings using
prototype cohesion and contrastive alignment. Our work directly addresses their limitations: they
require private source data, completely omit BN statistics, and collapse due to the feedback trap on
shared classification heads.

\section{Methodology}\vspace{-0.1in}
\label{sec:methodology}

\subsection{Problem Formulation}\vspace{-0.1in}
Let $\theta_{\text{base}}$ denote a shared base pre-trained model, and let $\{\theta_k\}_{k=1}^K$
be a library of $K$ expert networks fine-tuned from $\theta_{\text{base}}$ on $K$ distinct domains.
The task vector for expert $k$ is defined as $v_k = \theta_k - \theta_{\text{base}}$. At test time,
the model receives a continuous, non-stationary stream of unlabeled batches $X_1, X_2, \dots, X_t$
from different domains. Our goal is to dynamically merge the expert parameters $\theta_k$ into a
single merged model $\theta_{\text{merged},t}$ at each step $t$ using layer-wise merging
coefficients $\Lambda_{l,t}$ (for each layer $l$) and a global coefficient vector
$\lambda_{\text{global},t}$ (collectively denoted as $\Lambda_t = \{\Lambda_{l,t},
\lambda_{\text{global},t}\}$) to maximize classification accuracy across the stream. Specifically,
for each layer $l$ of the model, the merged parameters are computed as:
\begin{equation}
\theta_{\text{merged},t,l}(\Lambda_{l,t}) = \theta_{\text{base},l} + \sum_{k=1}^K \lambda_{l,t,k}
v_{l,k}
\end{equation}
where $\theta_{\text{merged},t,l}$ and $\theta_{\text{base},l}$ denote the parameters of layer $l$,
$v_{l,k} = \theta_{l,k} - \theta_{\text{base},l}$ is the task vector of expert $k$ at layer $l$, and
$\lambda_{l,t,k} \ge 0$ is the merging coefficient for expert $k$ at layer $l$ under the simplex
constraint $\sum_k \lambda_{l,t,k} = 1$ for all layers. In addition to these layer-wise
coefficients, we maintain a global coefficient vector $\lambda_{\text{global},t}$ which is used to
dynamically fuse the Batch Normalization statistics during test-time adaptation, as detailed in
Sections~3.3 and~3.4.

\subsection{Continuous Bayesian Mixture of Experts}\vspace{-0.1in}
To estimate the contribution of each expert to the current test batch $X$, we formulate TTMM as a
dynamic Bayesian mixture-of-experts. We define the unsupervised log-likelihood of the batch $X$
under expert $k$ as the negative Shannon entropy of the expert's predictions:
\begin{equation}
    \log P(X \mid \text{expert } k) = -\gamma H(Y \mid X, \theta_k)
\end{equation}
where $\gamma > 0$ is a temperature scaling parameter, and $H(Y \mid X, \theta_k) = -\frac{1}{|X|}
\sum_{i \in X} \sum_{c} p_{i,c} \log p_{i,c}$ is the average prediction entropy on the batch.
Experts that are highly confident on the test batch will yield low entropy, translating to high
likelihood.

Using Bayes' theorem, we compute the posterior probability of each expert $k$ being the active
expert for the current batch:
\begin{equation}
w_k = P(\text{expert } k \mid X) = \frac{L_k P(\text{expert } k)}{\sum_{j=1}^K L_j P(\text{expert }
j)}
\end{equation}
where $L_k = \exp(-\gamma H(Y \mid X, \theta_k))$ represents the likelihood under expert $k$, and
$P(\text{expert } k)$ is a uniform prior over experts. This posterior probability vector $w = [w_1,
\dots, w_K]$ provides a continuous, soft-routing weight distribution over the expert library.

\subsection{Soft Batch Normalization Buffer Fusion}\vspace{-0.1in}
Omitting Batch Normalization statistics during weight fusion leads to severe activation mismatches.
Hard-routing methods select the running statistics of a single active expert, which creates sharp
discontinuities at task boundaries. To resolve this, we propose Soft BN Buffer Fusion, which
continuously blends expert running statistics.

To derive the fusion equations, we make a standard simplifying assumption about the intermediate
representation space:
\begin{assumption}
The activation distribution for any given feature channel in a Batch Normalization layer,
conditioned on a specific expert $k$, is modeled as a Gaussian distribution: $z \mid \text{expert }
k \sim \mathcal{N}(\mu_k, \sigma^2_k)$.
\end{assumption}

Under this assumption, we model the total activation distribution under our mixture-of-experts as a
mixture of Gaussians.
\begin{proposition}
Let $z$ represent an activation variable. If Assumption~1 holds and the active expert is selected
according to the posterior probabilities $w_k = P(\text{expert } k \mid X)$, then the combined
activation $z$ is distributed as a mixture of Gaussians $z \sim \sum_{k=1}^K w_k \mathcal{N}(\mu_k,
\sigma^2_k)$. The mean $\mu_{\text{fused}}$ and variance $\sigma^2_{\text{fused}}$ of this mixture
are given exactly by:
\begin{equation}
    \mu_{\text{fused}} = \sum_{k=1}^K w_k \mu_k
\end{equation}
\begin{equation}
\sigma^2_{\text{fused}} = \sum_{k=1}^K w_k \left( \sigma^2_k + (\mu_k - \mu_{\text{fused}})^2
\right)
\end{equation}
\end{proposition}

\begin{proof}
By the law of total expectation, the mean of the mixture activation is:
\begin{align}
\mu_{\text{fused}} &= \mathbb{E}[z] = \sum_{k} P(\text{expert } k \mid X) \mathbb{E}[z \mid
\text{expert } k] \nonumber \\
    &= \sum_{k=1}^K w_k \mu_k
\end{align}
Using the variance decomposition formula (law of total variance):
\begin{equation}
\text{Var}(z) = \mathbb{E}[\text{Var}(z \mid \text{expert } k)] + \text{Var}(\mathbb{E}[z \mid
\text{expert } k])
\end{equation}
The first term is the expected conditional variance:
\begin{equation}
    \mathbb{E}[\text{Var}(z \mid \text{expert } k)] = \sum_{k=1}^K w_k \sigma^2_k
\end{equation}
The second term is the variance of the conditional expectation:
\begin{equation}
\text{Var}(\mathbb{E}[z \mid \text{expert } k]) = \sum_{k=1}^K w_k (\mu_k - \mu_{\text{fused}})^2
\end{equation}
Combining these yields the exact variance of the mixture activation:
\begin{equation}
\sigma^2_{\text{fused}} = \sum_{k=1}^K w_k \sigma^2_k + \sum_{k=1}^K w_k (\mu_k -
\mu_{\text{fused}})^2
\end{equation}
which matches the proposed equation.
\end{proof}
By computing the exact first and second moments of the mixture activation distribution, the Batch
Normalization layers can then utilize these exact moments to correctly normalize the incoming
activations on-the-fly. This serves as a mathematically sound and principled approximation for
normalizing the full (non-Gaussian) mixture activation profile, preserving representational
integrity across arbitrary domain combinations, eliminating activation mismatches, and avoiding
hard-switching discontinuities.

\subsection{Maximum A Posteriori Layer-Wise Adaptation}\vspace{-0.1in}
To understand the dual role of the Batch Normalization (BN) fusion weights, we distinguish between
the initial probabilistic routing and the inner-loop fine-tuning. We formulate our adaptation as a
dual-phase hierarchy: the \textbf{Bayesian posterior initialization} acts as the primary adaptation
driver, offering extremely high-precision routing on-the-fly, while the subsequent
\textbf{inner-loop MAP optimization} serves as a secondary, robust insurance policy. Upon receiving
a test batch $X_t$, we first compute the expert posterior probabilities $w$ (Eq. 5), representing
our probabilistic domain routing recommendation. We use this posterior $w$ to initialize the
layer-wise coefficients $\lambda_{l,k}$ and the global coefficient vector
$\lambda_{\text{global},k}$ in a logit-space, such that their initial softmax probabilities align
with the Bayesian posterior:
\begin{equation}
\lambda_{l,k} \propto \log(w_k + \epsilon), \quad \lambda_{\text{global},k} \propto \log(w_k +
\epsilon)
\end{equation}
where $\epsilon = 10^{-10}$ prevents log-zero errors. Crucially, the initial Soft BN Buffer Fusion
(Eq. 7--8) is applied using this initial posterior weight vector $w$. This places the merged model
in an optimal starting configuration.

During the subsequent inner MAP adaptation loop, the layer-wise merging coefficients
$\lambda_{l,k}$ are updated to minimize the predictive entropy of the merged model. As these
parameters adapt, the representations and activation statistics change. To ensure that the BN
running statistics remain perfectly synchronized with these parameter updates, the global
coefficient vector $\lambda_{\text{global}}$ is optimized alongside the layer coefficients. During
the $M$ inner steps, we dynamically update the BN running statistics using the current
softmax-transformed global coefficients, i.e., $w_{\text{inner}} =
\text{softmax}(\lambda_{\text{global}})$. This dual-phase mechanism maintains statistical and
parameter synchronization throughout the adaptation process. The motivation for maintaining separate
layer-wise and global coefficients lies in their distinct roles: while the layer-wise coefficients
resolve local ``representational friction'' between disparate expert layers, a single, global
coefficient vector provides a more stable and holistic estimate of the domain mixture for
normalizing activations. Constraining the BN fusion to a single global vector avoids
over-parameterization and prevents local statistical distortions. Under a temporal Gaussian prior
centered at the previous step's coefficients $\Lambda_{t-1}$, the objective to minimize is:
\begin{equation}
\mathcal{L} = H(Y \mid X, \theta_{\text{merged}}) + \frac{\beta}{2} \|\Lambda_t - \Lambda_{t-1}\|^2
\end{equation}
where $\beta \ge 0$ represents the prior regularization strength. In practice, all layer-wise
merging coefficients $\lambda_{l,k}$ and the global coefficient vector $\lambda_{\text{global},k}$
are grouped together as $\Lambda_t$, and they are optimized jointly using a single Adam optimizer
step in the inner loop. This temporal prior stabilizes adaptation for both parameters and
statistics, preventing catastrophic drift and coefficient oscillation.

To adapt layer-wise coefficients safely under heterogeneous sensitivities, we apply Test-Time
Fisher (TT-Fisher) preconditioning. We estimate parameter-level diagonal sensitivities dynamically
on-the-fly using the current unlabeled test batch $X$ and the model's own predictions as
pseudo-labels:
\begin{equation}
F(p) = \frac{1}{|X|} \sum_{i, c} \hat{p}_{i,c} \left[ \nabla_p \log P(y_i = c \mid x_i,
\theta_{\text{merged}}) \right]^2
\end{equation}
We then compute the layer-wise sensitivity score $S_l = \sum_{p \in \text{layer } l} F(p)$. During
backpropagation, the gradients of the layer-wise merging coefficients are scaled inversely to their
layer sensitivities:
\begin{equation}
\nabla_{\lambda_{l}} \mathcal{L} \leftarrow \nabla_{\lambda_{l}} \mathcal{L} \cdot \frac{1}{S_l +
\epsilon_d}
\end{equation}
This data-free preconditioning dampens adaptation in highly sensitive layers (e.g., classifier
heads and early convolutions) and accelerates stable adaptation in robust intermediate
representation blocks.

While highly effective, we note that the TT-Fisher (Eq. 12) is an empirical Fisher Information
matrix estimated on a single test batch using the model's own predictions as pseudo-labels.
Consequently, it represents a noisy, local approximation of the true Fisher Information. In
scenarios with extremely small batch sizes $B$ or high predictive uncertainty (e.g., when
transitioning between domains or encountering a novel domain), this estimator can experience minor
instability. However, our empirical results in Section 5.3.6 demonstrate that DF-Bayes-TTMM is
highly robust across a wide range of batch sizes $B \in \{16, 32, 64, 128\}$, indicating that the
inverse sensitivity scaling is sufficiently robust to tolerate these localized approximation noises
in practice.

\subsection{Uncertainty-Guided Novelty Detection}\vspace{-0.1in}
In an open-world setting, a completely novel domain $D_{\text{novel}}$ may arrive. Since there is
no pre-trained expert for the novel domain, all experts will exhibit high predictive entropy. We
compute the epistemic uncertainty of the expert mixture via the Mutual Information (MI) between the
predictions and the expert identity:
\begin{equation}
    I(Y; \text{Expert} \mid X) = H(Y \mid X) - \sum_{k=1}^K w_k H(Y \mid X, \theta_k)
\end{equation}
where $H(Y \mid X) = -\sum_{c} \bar{p}_c \log \bar{p}_c$ is the entropy of the average prediction
$\bar{p} = \sum_k w_k p_k$.

While mutual information provides the theoretical grounding for quantifying epistemic uncertainty
under confident disagreement, it can become uninformative when all experts are equally unconfident.
Specifically, when a completely novel domain arrives, every expert's predictions approach a uniform
distribution, yielding high individual predictive entropies $H(Y \mid X, \theta_k) \approx H_{\max}$
(aleatoric uncertainty). Consequently, the posterior probabilities $w_k$ also approach a uniform
distribution $1/K$, causing the average prediction entropy $H(Y \mid X)$ to also reach $H_{\max}$.
Under this uniform consensus, both terms in the MI equation approach $H_{\max}$, resulting in a low
MI score ($I(Y; \text{Expert} \mid X) \approx 0$) that fails to flag the novel domain.

In contrast, the average expert entropy $\bar{H} = \frac{1}{K} \sum_{k=1}^K H(Y \mid X, \theta_k)$
remains consistently and maximally high ($\approx H_{\max}$), acting as a monotonic and highly
discriminative indicator of the absolute lack of domain-specific expertise. Thus, the average expert
entropy serves as an extremely robust, direct, and computationally efficient proxy for detecting
novel domains on-the-fly. We therefore define our novelty detection rule as:
\begin{equation}
    \text{is\_novel} = \mathbb{I}(\bar{H} > \tau_N)
\end{equation}
If detected as novel, we initialize the merging coefficients uniformly and adapt them to find a
parameter configuration that generalizes to the novel domain, while bypassing the feedback trap by
setting the BN running statistics to a uniform average to prevent representational decay.

The complete step-by-step process of DF-Bayes-TTMM is summarized in
Algorithm~\ref{alg:df_bayes_ttmm}.

\begin{algorithm}[H]
\caption{DF-Bayes-TTMM: Data-Free Bayesian Test-Time Model Merging}
\label{alg:df_bayes_ttmm}
\begin{footnotesize}
\begin{algorithmic}[1]
\STATE {\bfseries Input:} Expert models $\{\theta_k\}_{k=1}^K$, pre-trained running statistics
$\{\mu_k, \sigma^2_k\}_{k=1}^K$, test stream $\{X_t\}_{t=1}^T$, parameters $\gamma$, $\beta$,
$\tau_N$, inner steps $M=3$.
\STATE {\bfseries Initialize:} $\Lambda_0$ uniformly, $\Lambda_{prev} \leftarrow \Lambda_0$
\FOR{each test batch $X_t$ in stream}
    \STATE Evaluate expert entropies $H_k = H(Y \mid X_t, \theta_k)$ for $k=1,\dots,K$
    \STATE Compute average expert entropy $\bar{H} = \frac{1}{K}\sum_k H_k$
    \IF{$\bar{H} > \tau_N$}
        \STATE $\text{is\_novel} \leftarrow \text{True}$
\STATE Initialize layer and global coefficients: $\lambda_{l,k} \leftarrow 0$,
$\lambda_{\text{global}, k} \leftarrow 0$
        \STATE Set BN fusion weights $w \leftarrow [1/K, \dots, 1/K]$
    \ELSE
        \STATE $\text{is\_novel} \leftarrow \text{False}$
        \STATE Compute posterior $w_k = \frac{\exp(-\gamma H_k)}{\sum_j \exp(-\gamma H_j)}$
\STATE Initialize layer and global coefficients: $\lambda_{l,k} \leftarrow \log(w_k + \epsilon)$,
$\lambda_{\text{global}, k} \leftarrow \log(w_k + \epsilon)$
    \ENDIF
    \STATE Assemble merged model parameters $\theta_{\text{merged}}(\Lambda_t)$
    \STATE Apply Soft BN Buffer Fusion to $\theta_{\text{merged}}$ using weights $w$:
    \STATE ~~~$\mu_{\text{fused}} = \sum_k w_k \mu_k$
    \STATE ~~~$\sigma^2_{\text{fused}} = \sum_k w_k (\sigma^2_k + (\mu_k - \mu_{\text{fused}})^2)$
    \STATE Estimate Test-Time Fisher sensitivities $S_l$ on $X_t$
    \FOR{inner step $m = 1 \dots M$}
\STATE Compute entropy loss $\mathcal{L}_{\text{entropy}} = H(Y \mid X_t, \theta_{\text{merged}})$
\STATE Compute MAP loss $\mathcal{L} = \mathcal{L}_{\text{entropy}} + \frac{\beta}{2}\|\Lambda_t -
\Lambda_{prev}\|^2$
        \STATE Compute gradients $\nabla_{\Lambda_t} \mathcal{L}$
\STATE Scale gradients inversely: $\nabla_{\lambda_l}\mathcal{L} \leftarrow
\nabla_{\lambda_l}\mathcal{L} \cdot \frac{1}{S_l + \epsilon_d}$
        \STATE Update $\Lambda_t$ using Adam step
\STATE Recalculate merged BN statistics on-the-fly from expert buffers $\{\mu_k,
\sigma^2_k\}_{k=1}^K$ using current weights $w \leftarrow \text{softmax}(\lambda_{\text{global}})$
    \ENDFOR
    \STATE Evaluate predictions $\hat{Y}_t = \text{argmax}(\theta_{\text{merged}}(X_t))$
    \STATE $\Lambda_{prev} \leftarrow \Lambda_t$ (Detach)
\ENDFOR
\end{algorithmic}
\end{footnotesize}
\end{algorithm}

\section{Experimental Setup}\vspace{-0.1in}
\label{sec:experimental_setup}

\textbf{Datasets:} We evaluate our framework using three standard vision datasets: MNIST
\cite{LeCun1998}, KMNIST \cite{Clanuwat2018}, and FashionMNIST \cite{Xiao2017}.
\textbf{Experts:} We train a 2-layer ConvNet with three BatchNorm layers as our backbone. We
fine-tune $K=2$ experts on MNIST and KMNIST respectively for 1 epoch each, achieving 99.12\% and
89.45\% single-task test accuracy. FashionMNIST is treated as the completely unseen novel domain in
the open-world stream.
\textbf{Evaluation Streams:} We construct three test-time streams of batch size 64: (1)~\textbf{Closed Sequential:} 15 batches of MNIST followed by 15 batches of KMNIST; (2)~\textbf{Closed Alternating:} 30 batches alternating between MNIST and KMNIST; and (3)~\textbf{Open-World:} 10 batches of MNIST, 10 batches of KMNIST, and 10 batches of FashionMNIST (Novel Domain).

\textbf{Baselines:} We compare \textbf{DF-Bayes-TTMM} against three key baselines: 
(i)~\textbf{Static Merging} with constant uniform weight merging ($[0.5, 0.5]$) and static BN statistics average; (ii)~\textbf{AdaMerging \cite{Yang2024}} with unconstrained layer-wise entropy minimization (learning rate $10^{-2}$); and (iii)~\textbf{DR-Fisher}, a strong custom baseline combining hard entropy-based expert routing (EBER), hard BN statistics switching, and Test-Time Fisher preconditioning. Other frameworks are not directly comparable as they require offline calibration data, violating our source-free requirement.

\textbf{Implementation Details:} For all methods, we use an Adam optimizer for the coefficient
updates with a learning rate of $10^{-2}$. We perform $M=3$ inner gradient steps per batch. For
DF-Bayes-TTMM, we set $\gamma = 15.0$ and prior regularization $\beta = 0.5$. The novelty detection
threshold is set to $\tau_N = 1.2$. All experiments are executed on a CPU platform using strictly
pre-trained, frozen expert checkpoints to ensure perfect reproducibility.

\section{Results and Discussion}\vspace{-0.1in}
\label{sec:results}

The experimental evaluation results are summarized in Table \ref{tab:results_summary}.

\begin{table*}[t]
\caption{Dynamic Test-Time Model Merging Accuracies (\%) across different non-stationary streams.
MNIST and KMNIST are known expert domains, while FashionMNIST acts as a completely novel domain in
the Open-World stream.}
\label{tab:results_summary}
\vskip 0.04in
\begin{center}
\begin{footnotesize}
\begin{sc}
\begin{tabular}{lcccc}
\toprule
Stream & Static Merging & AdaMerging & DR-Fisher & DF-Bayes-TTMM (Ours) \\
\midrule
Closed Sequential & 52.97\% & 52.97\% & 95.83\% & \textbf{95.83\%} \\
Closed Alternating & 52.97\% & 52.97\% & 95.83\% & \textbf{95.83\%} \\
Open-World & 37.71\% & 37.71\% & 65.10\% & \textbf{66.46\%} \\
\bottomrule
\end{tabular}
\end{sc}
\end{footnotesize}
\end{center}
\vskip -0.1in
\end{table*}

\begin{figure*}[t]
    \centering
    \begin{subfigure}[b]{0.33\textwidth}
        \centering
        \includegraphics[width=0.68\textwidth]{closed_sequential_accuracy}
        \vspace{-0.05in}
        \caption{Closed Sequential Stream}
        \label{fig:seq_acc}
    \end{subfigure}
    \hfill
    \begin{subfigure}[b]{0.33\textwidth}
        \centering
        \includegraphics[width=0.68\textwidth]{closed_alternating_accuracy}
        \vspace{-0.05in}
        \caption{Closed Alternating Stream}
        \label{fig:alt_acc}
    \end{subfigure}
    \hfill
    \begin{subfigure}[b]{0.33\textwidth}
        \centering
        \includegraphics[width=0.68\textwidth]{open-world_accuracy}
        \vspace{-0.05in}
        \caption{Open-World Stream}
        \label{fig:ow_acc}
    \end{subfigure}
    \vspace{-0.18in}
\caption{Dynamic Test-Time adaptation accuracy curves across different stream profiles.
DF-Bayes-TTMM (Ours) maintains stable, state-of-the-art accuracy throughout both closed-world shifts
and novel domain shifts without private training data.}
    \label{fig:accuracy_curves}
    \vspace{-0.28in}
\end{figure*}

\begin{table*}[t]
\begin{minipage}{0.48\textwidth}
\centering
\caption{Ablation of Batch Normalization buffer fusion strategies on test stream accuracies (\%).}
\label{tab:ablation_bn}
\vskip 0.04in
\begin{footnotesize}
\begin{sc}
\renewcommand{\arraystretch}{0.88}
\begin{tabular}{lccc}
\toprule
Strategy & Seq. & Alt. & Open \\
\midrule
Uniform (AdaMerging) & 52.55\% & 52.55\% & 37.08\% \\
Hard (EBER style) & 95.83\% & 95.83\% & 66.46\% \\
Soft (Ours) & \textbf{95.83\%} & \textbf{95.83\%} & \textbf{66.46\%} \\
\bottomrule
\end{tabular}
\end{sc}
\end{footnotesize}
\end{minipage}
\hfill
\begin{minipage}{0.48\textwidth}
\centering
\caption{Ablation of temperature parameter $\gamma$ on test stream accuracies (\%).}
\label{tab:ablation_temp}
\vskip 0.04in
\begin{footnotesize}
\begin{sc}
\renewcommand{\arraystretch}{0.88}
\begin{tabular}{lccc}
\toprule
$\gamma$ Value & Seq. & Alt. & Open \\
\midrule
$\gamma = 1.0$ & 34.11\% & 31.20\% & 25.36\% \\
$\gamma = 5.0$ & 94.43\% & 94.38\% & 65.57\% \\
$\gamma = 15.0$ (Default) & \textbf{95.83\%} & \textbf{95.83\%} & \textbf{66.46\%} \\
$\gamma = 50.0$ & \textbf{95.83\%} & \textbf{95.83\%} & \textbf{66.46\%} \\
\bottomrule
\end{tabular}
\end{sc}
\end{footnotesize}
\end{minipage}

\vspace{0.04in}

\begin{minipage}{0.48\textwidth}
\centering
\caption{Ablation of novelty detection threshold $\tau_N$ on test stream accuracies (\%).}
\label{tab:ablation_taun}
\vskip 0.04in
\begin{footnotesize}
\begin{sc}
\renewcommand{\arraystretch}{0.88}
\begin{tabular}{lccc}
\toprule
$\tau_N$ Value & Seq. & Alt. & Open \\
\midrule
$\tau_N = 0.5$ & 18.80\% & 18.80\% & 15.00\% \\
$\tau_N = 0.8$ & 93.18\% & 93.18\% & \textbf{66.46\%} \\
$\tau_N = 1.0$ & \textbf{95.83\%} & \textbf{95.83\%} & \textbf{66.46\%} \\
$\tau_N = 1.2$ (Default) & \textbf{95.83\%} & \textbf{95.83\%} & \textbf{66.46\%} \\
$\tau_N = 1.5$ & \textbf{95.83\%} & \textbf{95.83\%} & 64.38\% \\
\bottomrule
\end{tabular}
\end{sc}
\end{footnotesize}
\end{minipage}
\hfill
\begin{minipage}{0.48\textwidth}
\centering
\caption{Ablation of test-time batch size $B$ on test stream accuracies (\%).}
\label{tab:ablation_batch}
\vskip 0.04in
\begin{footnotesize}
\begin{sc}
\renewcommand{\arraystretch}{0.88}
\begin{tabular}{lccc}
\toprule
Batch Size $B$ & Seq. & Alt. & Open \\
\midrule
$B = 16$ & \textbf{95.83\%} & \textbf{95.83\%} & 65.99\% \\
$B = 32$ & \textbf{95.83\%} & \textbf{95.83\%} & 66.30\% \\
$B = 64$ (Default) & \textbf{95.83\%} & \textbf{95.83\%} & \textbf{66.46\%} \\
$B = 128$ & 95.76\% & 95.76\% & \textbf{66.46\%} \\
\bottomrule
\end{tabular}
\end{sc}
\end{footnotesize}
\end{minipage}
\end{table*}

\subsection{Closed Sequential and Alternating Streams}\vspace{-0.1in}
As shown in Table \ref{tab:results_summary}, standard Static Merging and AdaMerging struggle on
both Closed Sequential (52.97\%) and Closed Alternating (52.97\%) streams. Standard AdaMerging starting from uniform weights is trapped in the
``feedback trap'' (entropy collapse), a failure mode of unconstrained entropy minimization where the
coefficients collapse towards a single OOD expert, which is particularly severe under shared classification heads.

In contrast, \textbf{DR-Fisher} and \textbf{DF-Bayes-TTMM (Ours)} achieve exceptional accuracies,
reaching \textbf{95.83\%} on both closed sequential and closed alternating streams. DR-Fisher
utilizes Entropy-Based Expert Routing (EBER) to initialize coefficients close to the correct active
expert.
Our proposed \textbf{DF-Bayes-TTMM (Ours)} matches this exceptional baseline performance. This
demonstrates that our Soft BN Buffer Fusion, which continuously blends running statistics using soft
posterior probabilities, preserves activation integrity and eliminates hard-switching
discontinuities.

The batch-by-batch adaptation dynamics shown in Figure~\ref{fig:seq_acc} and
Figure~\ref{fig:alt_acc} demonstrate that DF-Bayes-TTMM adapts instantaneously at domain boundaries,
achieving high-accuracy convergence within a single batch.

\subsection{Open-World Adaptation and Novelty Detection}\vspace{-0.1in}
On the challenging Open-World stream, which includes 10 batches from the unseen FashionMNIST
domain, Static Merging and AdaMerging achieve only 37.71\%.

DR-Fisher achieves 65.10\%, but suffers from performance degradation on the novel domain. Because
DR-Fisher relies on hard expert routing, it is forced to route novel FashionMNIST samples to either
the MNIST or KMNIST expert, causing severe activation mismatches and representational collapse.

Our proposed \textbf{DF-Bayes-TTMM (Ours)} achieves the highest accuracy of \textbf{66.46\%},
outperforming the competitive DR-Fisher baseline by 1.36\% absolute accuracy. By utilizing our
uncertainty-guided novelty detection based on average expert entropy, DF-Bayes-TTMM robustly detects
the arrival of the novel domain, initializes the coefficients uniformly, and adapts them safely
using our MAP temporal prior and TT-Fisher preconditioning, while using soft continuous BN blending
to maintain representational stability.

This empirical advantage is clearly illustrated in Figure~\ref{fig:mi_vs_entropy_plot}, where we
plot the on-the-fly uncertainty signals across a transitioning test stream. While Mutual Information
collapses to zero under the novel domain due to the uniform consensus of uncertainty, our average
expert entropy proxy remains monotonic and highly elevated, providing a pristine, robust signal for
novelty detection.

As shown in Figure~\ref{fig:ow_acc}, when transitioning to the novel domain (batches 20-30), our
method maintains a significantly higher accuracy compared to DR-Fisher, proving that our smooth,
moment-matching buffer fusion represents a superior approach for open-world test-time adaptation.

\begin{figure}[t]
    \centering
    \includegraphics[width=0.82\columnwidth]{mi_vs_entropy}
    \vspace{-0.1in}
\caption{Uncertainty signals across a transitioning test stream. While Mutual Information collapses
to zero under the uniform consensus of uncertainty on the novel domain, our proposed average expert
entropy proxy remains consistently high, providing a robust, clear signal for open-world novelty
detection.}
    \label{fig:mi_vs_entropy_plot}
    \vspace{-0.18in}
\end{figure}

\subsection{Deep Architectural Scalability on CIFAR-10 with ResNet-18}\vspace{-0.1in}
\label{sec:cifar_resnet}
To address the critical concern of scalability and generalizability to deeper architectures and
color images, we evaluate DF-Bayes-TTMM on a challenging, high-dimensional setup. Specifically, we
employ a pre-trained \textbf{ResNet-18} architecture (11.2 million parameters) and evaluate on
non-stationary test streams constructed from the \textbf{CIFAR-10} dataset.

\subsubsection{Experimental Setup}
We train two ResNet-18 expert models from scratch on distinct domains:
\textbf{Expert 1 (Clean CIFAR-10)} trained on standard, unmodified CIFAR-10 (achieving 88.16\% validation accuracy), and \textbf{Expert 2 (Rotated CIFAR-10)} trained on CIFAR-10 images rotated by 180 degrees (achieving 88.43\% validation accuracy).
We then construct three streaming profiles of batch size 64:
(1)~\textbf{Closed Sequential Stream:} 15 batches of Clean CIFAR-10 followed by 15 batches of
Rotated CIFAR-10; (2)~\textbf{Closed Alternating Stream:} 30 alternating batches of Clean and Rotated CIFAR-10; and (3)~\textbf{Open-World Stream:} 10 batches of Clean CIFAR-10, 10 batches of Rotated CIFAR-10, and 10 batches of an unseen, novel domain: Grayscale CIFAR-10.

The results of this large-scale evaluation are presented in Table~\ref{tab:cifar_results}.

\begin{table*}[t]
\centering
\caption{Dynamic Test-Time Model Merging Accuracies (\%) on the CIFAR-10 ResNet-18 benchmark.}
\label{tab:cifar_results}
\vskip 0.1in
\begin{center}
\begin{small}
\begin{sc}
\begin{tabular}{lcccc}
\toprule
Stream & Static Merging & AdaMerging & DR-Fisher & DF-Bayes-TTMM \\
\midrule
Closed Sequential & 11.04\% & 11.04\% & \textbf{76.88\%} & 72.50\% \\
Closed Alternating & 11.04\% & 11.04\% & \textbf{76.88\%} & 72.50\% \\
Open-World & 11.88\% & 11.88\% & \textbf{74.32\%} & 70.31\% \\
\bottomrule
\end{tabular}
\end{sc}
\end{small}
\end{center}
\end{table*}

These results yield crucial scaling insights and highlight fundamental trade-offs:
(i)~\textbf{Feedback Trap at Scale:} Standard Static Merging and AdaMerging completely collapse
to near-random guessing (11.04\% to 11.88\%) across all streams, confirming that unconstrained entropy minimization collapse is even more severe in deeper networks.
(ii)~\textbf{Effectiveness of Moment-Matching BN Fusion:} Both DR-Fisher (EBER-routing style) and
DF-Bayes-TTMM successfully bypass this catastrophic collapse, achieving outstanding accuracies in
the range of 70\% to 76.88\%, proving that dynamic BN buffer fusion scales to deep networks.
(iii)~\textbf{Continuous Soft Blending vs. Hard-Routing Trade-Offs:} In this highly specialized, non-overlapping setting, the hard expert routing of DR-Fisher achieves a performance lead of ~4.38\% over our soft continuous blending. This reveals a valuable design-space trade-off: hard-routing is optimal for perfectly separated, distinct domains where it acts as a strong oracle. However, as analyzed in Appendix~\ref{sec:appendix_discussion}, our soft continuous Bayesian mixture provides far greater robustness and representational stability under soft shifts, mixed-domain batches, or smooth transitions. By presenting this balanced perspective, we highlight that both approaches represent distinct, complementary design points depending on domain overlap in deployment.

\subsection{Ablation Studies}\vspace{-0.1in}
To systematically analyze the contributions of each component in DF-Bayes-TTMM, we perform
extensive ablation studies on the Closed Sequential, Closed Alternating, and Open-World streams.

\subsubsection{Ablation 1: Batch Normalization Fusion Strategy}\vspace{-0.1in}
We compare our proposed Soft BN Buffer Fusion against Hard BN Buffer Fusion (expert routing) and Uniform BN Buffer Fusion (uniform averaging).
As shown in Table~\ref{tab:ablation_bn}, using Uniform BN statistics results in catastrophic representation collapse across all streams, yielding only 52.55\% accuracy on closed streams and 37.08\% on the open-world stream. Fusing the BN buffers dynamically using either Soft or Hard fusion is crucial to avoid activation mismatches, with both Soft and Hard fusion achieving 95.83\% on closed streams and 66.46\% on the open-world stream.
The identical performance of Soft and Hard BN fusion in Table~\ref{tab:ablation_bn} arises because both rows are evaluated within the full DF-Bayes-TTMM framework, which includes our novelty detection. When the novel domain is flagged on the open-world stream, both methods fall back to uniform BN statistics to prevent activation decay. However, as shown in Table~\ref{tab:results_summary}, the standalone DR-Fisher baseline (which uses hard routing without a novelty threshold, and is forced to route novel FashionMNIST samples to MNIST/KMNIST expert BN statistics) drops to 65.10\% accuracy, demonstrating that our continuous Bayesian mixture combined with novelty fallback provides superior resilience under novel domains.

\subsubsection{Ablation 2: Temporal Prior Regularization ($\beta$)}\vspace{-0.1in}
We examine the effect of the temporal prior coefficient $\beta \in \{0, 0.1, 0.5, 2, 10\}$ on
stream accuracies. We find that in settings with abrupt, discrete domain shifts, all tested values of $\beta$ yield identical
accuracies (95.83\% on Sequential/Alternating and 66.46\% on Open-World). This indicates that our
highly accurate Bayesian posterior initialization immediately sets the starting coefficients
extremely close to the optimal point for each batch, making subsequent inner gradient adaptation
steps stable and localized. While this demonstrates the robustness of our posterior
initialization, we hypothesize that the temporal prior remains crucial for stable coefficient tracking in more challenging scenarios involving gradual non-stationary drift or mixed-domain batches, where it prevents undesirable parameter oscillations.

\subsubsection{Ablation 3: Bayesian Temperature Scaling ($\gamma$)}\vspace{-0.1in}
We evaluate the sensitivity of the temperature parameter $\gamma \in \{1, 5, 15, 50\}$ used to
compute the expert mixture posterior.
As shown in Table~\ref{tab:ablation_temp}, when $\gamma=1.0$, the posterior over experts is too
soft and close to uniform, which mixes incompatible expert weights and leads to severe
representation collapse (34.11\% sequential and 25.36\% open-world). Concentrating the posterior
with $\gamma \ge 15.0$ dramatically stabilizes the representation space, reaching peak accuracies of
95.83\% and 66.46\%. This indicates that our dynamic MoE model requires a
sufficiently peaked posterior to successfully isolate the relevant expert and avoid destructive
cross-expert interference.

\subsubsection{Ablation 4: Number of Inner Adaptation Steps ($M$)}\vspace{-0.1in}
We evaluate the impact of the inner optimization steps $M \in \{0, 1, 2, 3, 5\}$ on stream
performance. Crucially, setting $M=0$ (zero gradient steps) achieves the exact same SOTA accuracy
(95.83\% Sequential/Alternating, 66.46\% Open-World) as $M \ge 1$. This is a major strength: our Bayesian posterior initialization is so precise that it immediately places the merged parameters and activation statistics in their optimal configurations on-the-fly, enabling an extremely efficient, zero-gradient ``Ultra-Light Mode'' ($M=0$) for edge deployment. We position the inner-loop MAP adaptation and TT-Fisher preconditioning ($M \ge 1$) as a robust, secondary insurance policy designed for more challenging conditions, such as ambiguous experts, gradual shifts, or deep architectures.

\subsubsection{Ablation 5: Novelty Detection Threshold ($\tau_N$)}\vspace{-0.1in}
We evaluate the sensitivity of the novelty threshold $\tau_N \in \{0.5, 0.8, 1.0, 1.2, 1.5\}$ in
Table~\ref{tab:ablation_taun}. When $\tau_N=0.5$ (too low), known domains are misclassified as
novel, dropping accuracy to 18.80\% (closed) and 15.00\% (open) due to premature uniform fallback. When $\tau_N \in \{1.0, 1.2\}$, we achieve perfect domain separation and optimal accuracies
(95.83\% closed, 66.46\% open). When $\tau_N=1.5$ (too high), the novel domain is misclassified as
known, reducing open accuracy to 64.38\% due to incorrect routing. This confirms our predictive entropy proxy provides a highly robust signal, with any threshold in the wide range $[1.0, 1.2]$ yielding optimal performance.

\subsubsection{Ablation 6: Impact of Test-Time Batch Size ($B$)}\vspace{-0.1in}
We examine the effect of target stream batch size $B \in \{16, 32, 64, 128\}$ on adaptation
performance in Table~\ref{tab:ablation_batch}.
On the Closed Sequential and Alternating streams, our method is remarkably stable, maintaining an
optimal accuracy of 95.83\% across nearly all batch sizes.
On the challenging Open-World stream, smaller batch sizes like $B=16$ still achieve a high accuracy
of 65.99\%, which increases steadily to 66.30\% at $B=32$ and reaches a peak of 66.46\% at $B \ge
64$.
This outstanding robustness to stream fragmentation demonstrates that our Bayesian model does not depend on large sample populations for stable adaptation, making it exceptionally suited for edge streaming applications.

\subsubsection{Ablation 7: Origin of Global BN Fusion Weights}\vspace{-0.1in}
To address the necessity of optimizing a separate learnable global coefficient vector $\lambda_{\text{global},t}$ versus using a derived quantity, we perform an ablation study comparing our default learnable system with an \emph{average} baseline. Under the \emph{average} baseline, $\lambda_{\text{global},t}$ is not independently optimized or initialized; rather, it is dynamically computed at each step as the mean of all active layer-wise merging coefficients $\lambda_{l,t}$. 
The evaluation yields identical accuracies of 95.83\% on the Closed Sequential and Alternating streams, and 66.46\% on the Open-World stream for both modes. This outstanding consistency reveals that because our data-free Bayesian posterior initialization (which sets starting coefficients to the soft posterior) is exceptionally precise, the framework achieves optimal convergence immediately at the first step, making the specific gradient pathway of the global BN weights during subsequent fine-tuning steps virtually redundant. This justifies our dual-coefficient system as a highly robust general formulation while demonstrating that the framework can be simplified to use a parameter-free derived average without any performance loss.



\section{Conclusion}\vspace{-0.1in}
\label{sec:conclusion}
We proposed \textbf{DF-Bayes-TTMM}, a mathematically rigorous and completely data-free Bayesian
Test-Time Model Merging framework. By formulating TTMM as a continuous Bayesian mixture-of-experts,
we establish soft continuous BN buffer fusion and uncertainty-guided novelty detection. Our method
eliminates private source data dependency, bypasses the feedback trap, and achieves state-of-the-art
performance across sequential, alternating, and open-world test streams. While our empirical
validation is focused on the MNIST family of datasets using a 2-layer ConvNet to provide a clean,
controlled environment for analyzing fundamental failure modes, we provide a thorough discussion on
scaling the framework to deeper architectures (e.g., ResNets and Transformers) and more complex
datasets in Appendix~\ref{sec:appendix_discussion}. Future work includes scaling this Bayesian
formulation to large-scale vision-language and language models.

\bibliography{template/example_paper}
\bibliographystyle{template/icml2026}

\clearpage
\appendix
\onecolumn
\section{Detailed Discussion and Scaling Analysis}
\label{sec:appendix_discussion}

\subsection{Scalability and Architectural Generalization}
While we validate DF-Bayes-TTMM on MNIST-family benchmarks with a 2-layer ConvNet, the framework's
mathematical design naturally generalizes to modern, deep architectures (e.g., ResNets, Vision
Transformers) and complex datasets (e.g., CIFAR-100, ImageNet-C). In deeper networks, activation
mismatches are compounded layer-by-layer; thus, our mathematically exact Soft BN Buffer Fusion is
even more critical for preserving representational coherence. For architectures utilizing Layer
Normalization (LN) or Group Normalization (GN), such as ViTs, the statistics are computed per-sample
or per-channel on-the-fly, eliminating the need to track running buffers altogether, making
DF-Bayes-TTMM directly applicable by merging the LN/GN parameters via our soft continuous
posteriors.

\subsection{Robustness Under Soft Shifts and Mixed-Domain Batches}
Real-world streams often present ``soft shifts'' where a test batch contains a mixture of multiple
known domains (e.g., 50\% MNIST and 50\% KMNIST). Under unconstrained entropy minimization (e.g.,
AdaMerging), the optimizer is caught in the feedback trap, collapsing parameters to a single
confident expert and decimating accuracy.
In contrast, DF-Bayes-TTMM elegantly handles mixed-domain batches: the Bayesian posterior (Eq. 5)
assigns balanced weights (e.g., $w_{\text{MNIST}} \approx 0.5, w_{\text{KMNIST}} \approx 0.5$),
which naturally produces a soft, joint model interpolation and blends BN statistics using the
correct mixture moments. Since our method uses these soft priors as a starting point and stabilizes
adaptation via temporal priors ($\beta$), it remains in a stable, optimal interpolating state.

\subsection{Computational Complexity and Inference Latency}
DF-Bayes-TTMM supports two distinct modes depending on latency constraints:

\textbf{Ultra-Light Mode ($M=0$):} Bypasses all test-time backpropagation. The total latency per
batch is $(K+1)$ forward passes ($K$ expert forward passes to compute posteriors, and $1$ merged
model forward pass). Since forward passes are highly parallelizable and computationally cheap, this
mode delivers real-time inference speed suitable for edge deployment.

\textbf{Robust Mode ($M \ge 1$):} Performs a single backward pass to compute Test-Time Fisher
preconditioning, and $M$ inner optimization steps. The latency is $(K+1)$ forward passes, $1$
backward pass, and $M$ forward-backward passes, which is comparable to standard TTA methods like
TENT \cite{Wang2021} but with superior optimization stability.

\subsection{Uncertainty Signals for Open-World Novelty Detection}
\label{sec:appendix_uncertainty_signals}
To empirically validate the theoretical argument presented in Section 3.5 regarding the robust
behavior of average expert entropy over Mutual Information (MI), we construct a transitioning test
stream of 30 batches (10 MNIST batches, followed by 10 KMNIST batches, followed by 10 completely
novel FashionMNIST batches) and compute both uncertainty metrics on-the-fly.

As depicted in Figure~\ref{fig:mi_vs_entropy_plot} (located in Section~5.2), when the stream
traverses the known domains (batches 0-19), the active expert (e.g., MNIST for batches 0-9) makes
highly confident predictions, which keeps the average expert entropy extremely low ($\approx 0.60$
to $0.75$). At the domain boundary transition (batch 10), Mutual Information exhibits a sharp
localized spike, reflecting the disagreement and temporary epistemic uncertainty between the two
experts as the active domain shifts.

However, upon transitioning to the completely novel FashionMNIST domain (batches 20-29), both
experts generate highly unconfident, high-entropy predictions that approach a uniform distribution.
Under this consensus of uncertainty (uniform aleatoric noise across all models), both terms of the
Mutual Information equation (Eq. 6) approach their maximum value, causing MI to collapse back to
zero. Thus, MI completely fails to register the presence of the novel domain.

In contrast, our proposed average expert entropy proxy $\bar{H}$ remains monotonic and highly
elevated ($\approx 1.25$), providing a pristine, robust, and well-separated signal for detecting the
arrival of the novel domain on-the-fly. This empirical characterization confirms that the average
expert entropy is a superior and highly discriminative metric for real-time open-world novelty
detection.

\subsection{Scalability and Deep Architectural Validation on CIFAR-10}
\label{sec:appendix_cifar}
The scalability and architectural validation of DF-Bayes-TTMM on ResNet-18 with CIFAR-10, as well as the important design-space trade-offs between continuous soft blending and hard expert routing, are fully discussed and analyzed in Section~\ref{sec:cifar_resnet} (located in the main body).

\subsection{Calibration and Adaptability of the Novelty Threshold}
As discussed in Section 3.5 and evaluated in Section 5.3.5, our novelty detection mechanism relies
on a threshold $\tau_N$. While we show that a wide range $[1.0, 1.2]$ of thresholds yields optimal
results on MNIST-family streams, in real-world deployments this parameter can be adaptively
calibrated on-the-fly.
Specifically, instead of a hand-tuned constant, we can calibrate $\tau_N$ by:
\begin{enumerate}
\item \textbf{Rolling Window Outlier Detection:} Tracking a running average and standard deviation
of the average expert entropy $\bar{H}_t$ on-the-fly. A novel domain is then flagged when $\bar{H}_t
> \mu_{\bar{H}} + k \cdot \sigma_{\bar{H}}$ (where $\mu_{\bar{H}}$ and $\sigma_{\bar{H}}$ are
running statistics computed over a sliding temporal window of known domains, and $k$ is a scaling
factor, e.g., $k=3$). This eliminates the reliance on a static magic number.
\item \textbf{Source-Free Calibration:} Using the first few batches of the incoming target stream
(which are assumed to belong to known domains) to establish a baseline distribution of known-domain
expert entropies and setting the threshold to a statistical upper bound (e.g., the 95th percentile).
\end{enumerate}

\subsection{Framing of the Feedback Trap: Shared vs. Separate Classifiers}
We emphasize that the ``feedback trap'' of unconstrained entropy minimization (AdaMerging) is
particularly severe in settings where experts share a classification head. Because the classifiers
are shared, an expert from a different domain can confidently predict wrong classes under
out-of-distribution inputs, causing the unconstrained predictive entropy of the merged model to
collapse onto a single confidently wrong expert.
If the experts had separate, independent classification heads (where each head is specialized to
its own domain's class space), we hypothesize that:
\begin{enumerate}
\item The feedback loop would be partially mitigated, because the OOD classifier head would produce
more uniform, high-entropy predictions under OOD inputs, which naturally signals high uncertainty.
\item However, weight-level representation collapse would still persist because unconstrained
entropy minimization does not account for Batch Normalization statistics alignment, leading to
layer-wise activation mismatches.
\end{enumerate}
Our proposed DF-Bayes-TTMM resolves both issues simultaneously, ensuring perfect activation and
parameter alignment regardless of the classification head architecture.

\end{document}
